\documentclass[11pt, a4paper]{article}
\usepackage[a4paper, top=2.5cm, bottom=2.5cm, left=2cm, right=2cm]{geometry}
\usepackage{fontspec}
\usepackage[english, bidi=basic, provide=*]{babel}
\babelprovide[import, onchar=ids fonts]{english}
\babelfont{rm}{TeX Gyre Termes}
\babelfont{sf}{TeX Gyre Heros}
\babelfont{tt}{TeX Gyre Cursor}
\usepackage{enumitem}
\usepackage{amsmath}

\title{\textbf{Chapter 5: The Structure and Function of Large Biological Molecules \\ Exam Notes}}
\author{}
\date{}

\begin{document}

\maketitle

\section*{5.1 Macromolecules are Polymers, Built from Monomers}
Large biological molecules form a diverse array of complex structures that are built from simpler, repeating units.
\begin{itemize}
    \item \textbf{Macromolecules:} Giant molecules formed by the joining of smaller molecules. The three truly polymeric macromolecules are Carbohydrates, Proteins, and Nucleic Acids. (Lipids are large but are not strictly polymers).
    \item \textbf{Polymers and Monomers:} A \textit{polymer} is a long molecule consisting of many similar or identical building blocks linked by covalent bonds. The repeating units are called \textit{monomers}.
    \item \textbf{Dehydration Reaction (Synthesis):} Monomers are connected by a reaction in which two molecules are covalently bonded to each other, with the \textbf{loss of a water molecule}. One monomer provides a hydroxyl group ($-OH$) and the other provides a hydrogen ($-H$). This process requires energy and is facilitated by enzymes.
    \item \textbf{Hydrolysis (Breakdown):} The reverse of dehydration. Polymers are disassembled to monomers by adding a water molecule, breaking the covalent bond (e.g., digestion in the human body).
\end{itemize}

\section*{5.2 Carbohydrates: Fuel and Building Material}
Carbohydrates include sugars and the polymers of sugars.
\begin{itemize}
    \item \textbf{Monosaccharides:} The simplest carbohydrates (simple sugars), generally with molecular formulas that are multiples of $CH_2O$. 
    \begin{itemize}
        \item \textit{Glucose ($C_6H_{12}O_6$)} is the most common. 
        \item Classified by the location of the carbonyl group (aldose vs. ketose) and the size of the carbon skeleton (e.g., hexoses, pentoses).
        \item Form ring structures in aqueous solutions.
    \end{itemize}
    \item \textbf{Disaccharides:} Consist of two monosaccharides joined by a \textbf{glycosidic linkage} (a covalent bond formed via a dehydration reaction). Examples: Sucrose (glucose + fructose), Lactose, Maltose.
    \item \textbf{Polysaccharides (Storage):}
    \begin{itemize}
        \item \textit{Starch:} A polymer of glucose monomers used by \textbf{plants} to stockpile surplus glucose. Synthesized with $\alpha$ (alpha) glucose rings, causing helical structures.
        \item \textit{Glycogen:} A highly branched polymer of glucose used by \textbf{animals}. Stored mainly in liver and muscle cells. Depletes in about a day unless replenished.
    \end{itemize}
    \item \textbf{Polysaccharides (Structural):}
    \begin{itemize}
        \item \textit{Cellulose:} A major component of tough plant cell walls. It is a polymer of $\beta$ (beta) glucose, forming straight, unbranched chains. Hydrogen bonds form between parallel cellulose molecules, creating strong microfibrils. Most animals (including humans) cannot digest cellulose.
        \item \textit{Chitin:} The carbohydrate used by arthropods (insects, crustaceans) to build exoskeletons, and by fungi in their cell walls.
    \end{itemize}
\end{itemize}

\section*{5.3 Lipids: A Diverse Group of Hydrophobic Molecules}
Lipids are grouped together because they mix poorly, if at all, with water (they are highly hydrophobic due to consisting mostly of hydrocarbon regions).
\begin{itemize}
    \item \textbf{Fats (Triglycerides):} Constructed from glycerol and three fatty acids joined by an \textbf{ester linkage}. Used primarily for long-term energy storage, insulation, and cushioning.
    \begin{itemize}
        \item \textit{Saturated Fatty Acids:} Have no double bonds between carbon atoms. The chain is straight, allowing molecules to pack tightly. Solid at room temp (e.g., butter, animal fats).
        \item \textit{Unsaturated Fatty Acids:} Have one or more double bonds (nearly always \textit{cis} double bonds), which create a "kink" in the chain. Cannot pack tightly. Liquid at room temp (e.g., olive oil).
    \end{itemize}
    \item \textbf{Phospholipids:} Essential for cells because they make up cell membranes. Consist of a glycerol attached to two fatty acids and a phosphate group.
    \begin{itemize}
        \item \textit{Amphipathic:} They have a hydrophilic (polar) head and two hydrophobic (nonpolar) tails.
        \item In water, they self-assemble into double-layered structures called \textbf{bilayers}, shielding the hydrophobic portions from water.
    \end{itemize}
    \item \textbf{Steroids:} Lipids characterized by a carbon skeleton consisting of four fused rings.
    \begin{itemize}
        \item \textit{Cholesterol:} A crucial steroid in animals; a component of cell membranes and a precursor from which other steroids (like sex hormones) are synthesized.
    \end{itemize}
\end{itemize}

\section*{5.4 Proteins: A Diversity of Structures and Functions}
Proteins account for more than 50\% of the dry mass of most cells and are instrumental in almost everything organisms do (catalysis, defense, transport, cellular communication, movement).
\begin{itemize}
    \item \textbf{Amino Acids:} The monomers of proteins. They possess both an amino group and a carboxyl group, an $\alpha$ (alpha) carbon, a hydrogen atom, and a variable side chain (\textbf{R group}).
    \item \textbf{R Groups:} Determine the unique characteristics of a particular amino acid (nonpolar/hydrophobic, polar/hydrophilic, acidic/negative, basic/positive). There are 20 standard amino acids.
    \item \textbf{Polypeptides:} Amino acids are linked by \textbf{peptide bonds} (formed by dehydration reactions). A polypeptide has a repeating backbone with free amino and carboxyl ends (N-terminus and C-terminus).
    \item \textbf{Four Levels of Protein Structure:}
    \begin{enumerate}
        \item \textit{Primary Structure:} The linear sequence of amino acids, determined by inherited genetic information.
        \item \textit{Secondary Structure:} Regions stabilized by hydrogen bonds between atoms of the polypeptide backbone. Common structures are the $\alpha$ helix (coils) and $\beta$ pleated sheet (folds).
        \item \textit{Tertiary Structure:} The overall 3D shape of a polypeptide resulting from interactions between the \textbf{R groups} (hydrophobic interactions, van der Waals interactions, hydrogen bonds, ionic bonds, and strong covalent \textbf{disulfide bridges}).
        \item \textit{Quaternary Structure:} Arises when a protein consists of two or more polypeptide chains aggregated into one functional macromolecule (e.g., Collagen, Hemoglobin).
    \end{enumerate}
    \item \textbf{Denaturation:} The loss of a protein's native shape due to changes in pH, salt concentration, or temperature. A denatured protein is biologically inactive.
    \item \textbf{Chaperonins:} Protein molecules that assist the proper folding of other proteins by keeping the new polypeptide segregated from "bad influences" in the cytoplasm.
\end{itemize}

\section*{5.5 Nucleic Acids: Storing and Transmitting Information}
Nucleic acids store, transmit, and help express hereditary information. The amino acid sequence of a polypeptide is programmed by a unit of inheritance known as a gene, consisting of DNA.
\begin{itemize}
    \item \textbf{Two Types:}
    \begin{itemize}
        \item \textit{Deoxyribonucleic acid (DNA):} Provides directions for its own replication. Directs RNA synthesis.
        \item \textit{Ribonucleic acid (RNA):} Controls protein synthesis (Gene expression flow: $DNA \rightarrow RNA \rightarrow Protein$).
    \end{itemize}
    \item \textbf{Nucleotides:} The monomers of nucleic acids. Consist of three parts:
    \begin{enumerate}
        \item A nitrogenous base.
        \item A five-carbon sugar (pentose).
        \item One or more phosphate groups.
    \end{enumerate}
    \item \textbf{Nitrogenous Bases:}
    \begin{itemize}
        \item \textit{Pyrimidines:} Have a single six-membered ring. Cytosine (C), Thymine (T, only in DNA), Uracil (U, only in RNA).
        \item \textit{Purines:} Have a six-membered ring fused to a five-membered ring. Adenine (A), Guanine (G).
    \end{itemize}
    \item \textbf{Sugars:} Deoxyribose in DNA (lacks an oxygen atom on the second carbon) and Ribose in RNA.
    \item \textbf{Polynucleotides:} Nucleotides are linked by \textbf{phosphodiester linkages} creating a sugar-phosphate backbone. They have built-in directionality (5' end to 3' end).
    \item \textbf{DNA Double Helix:} DNA molecules have two polynucleotide strands spiraling around an imaginary axis. The two backbones run in opposite 5' $\rightarrow$ 3' directions (\textbf{antiparallel}). Base pairing is strict: Adenine pairs with Thymine ($A-T$), and Guanine pairs with Cytosine ($G-C$).
\end{itemize}

\section*{5.6 Genomics and Proteomics}
\begin{itemize}
    \item \textbf{Genomics:} The approach of analyzing large sets of genes or even comparing whole genomes of different species.
    \item \textbf{Proteomics:} The analysis of large sets of proteins, including their sequences.
    \item \textbf{Bioinformatics:} The use of computer software and computational tools to handle and analyze the massive data sets resulting from sequencing genomes.
    \item \textbf{Evolutionary Significance:} Linear sequences of DNA are passed from parents to offspring. Species that are more closely related share a greater proportion of their DNA and protein sequences.
\end{itemize}

\end{document}
